\chapter{Life}

\section{Definition of Life}

Life appears to be difficult to define.

The first obvious problem is those borderline cases. Entities such as viruses, prions, and sterile organisms
satisfy some criteria for life but not others. Any strict definition tends to exclude things many
scientists consider ``life'' or include things they do not.

Recent developments in computers and AI raise another problem, namely substrate independence.
We can build digital replications of living cells in computers. Can those digital replicas be
considered ``life'' as well?

Observer bias is also a problem. Our definitions are based on life on Earth, so it is unclear whether
they would capture fundamentally different forms of life elsewhere.


\section{Evolution}

There is a science fiction book, \emph{The Black Cloud} \cite{hoyle1957blackcloud},
in which an intelligent cloud arrives and proceeds to cause all sorts
of trouble. The reason the book is science fiction rather than science,
and why intelligent clouds cannot exist in real life, is that there are
no known laws of physics on which such a cloud could plausibly be based.

The more intelligent and complex a system is, the smaller the
probability that it could simply appear spontaneously. For a giant
intelligent cloud, the probability should be practically zero.

The only theory we currently have that explains the existence of
truly intelligent systems is evolution. For evolution to work there
must be many candidates and a mechanism of natural selection capable
of eliminating the less successful ones. In the case of intelligent
clouds, there would have to be lots of clouds, and natural selection
would need to eliminate the weak and disorganized ones while favoring
those capable of maintaining structure. Only then could intelligent
clouds gradually develop.

But how would evolution work for a cloud whose behavior is governed
by the Navier--Stokes equations of fluid dynamics? How could such a
cloud keep its information in order? What would happen if a storm
passed by and blew the cloud apart? Even relatively small disturbances to a human brain can cause severe
damage. A violent disturbance would correspond to putting a brain into
a kitchen blender and switching it on. It is not difficult to see why
the brain would not think very clearly afterward.

In principle one could imagine some unknown form of matter with
properties suitable for maintaining information in a gaseous state.
However, astronomical observations give us little support for this
idea. When we analyze the spectrum of light coming from anywhere in
the universe—even from the most distant galaxies—we see essentially
the same electromagnetic fingerprints. This tells us that they are
made of the same raw materials as our own home sweet Milky Way.

An intelligent cloud would need to know its boundaries and maintain
a stable internal structure. A gas cloud has neither. It would not
remain organized long enough to evolve intelligence.

Intelligence requires stable structures capable of storing and
protecting information. Turbulent gases are exceptionally poor at
doing this.

While Charles Darwin's theory of evolution has been subject to heavy
religion-saturated debate in the past, believing in it is no longer a
question of faith. Modern science has revealed the genetic code in DNA
and confirmed the theory for all its essential parts. It appears that
all life on Earth shares a common ancestor.


\section{Alien Life}

Life could be far more common than generally believed and might not necessarily be carbon-based.
However, the challenge is that Earth appears to be an ideal habitat for diverse life forms.
If life emerged easily, shouldn't we observe "exotic" organisms—perhaps non-DNA-based—coexisting with us? Instead, we find only one biochemical lineage, all rooted in DNA and evolution.
Furthermore, our exploration of the solar system and our extensive monitoring of radio signals (SETI) have yielded no traces of life, intelligent or otherwise.

Perhaps life exists on a different temporal scale. If their biological processes run significantly faster or slower than ours, we might perceive them as either stationary matter or a blur too rapid to recognize as sentient.
However, while time-scale differences are a great sci-fi concept, biology is bound by the laws of thermodynamics and chemistry.
Chemical reactions (the basis of life) happen at specific rates dictated by temperature and molecular stability.
A creature "moving too fast" would likely burn up from the heat of its own metabolism; one "moving too slow" might not be able to gather enough energy to maintain its structure against entropy.

Maybe exotic life could be here, but we simply aren't looking for it correctly.
Most of our tools (PCR, DNA sequencing) are designed specifically to find DNA.
If a "non-DNA" microbe existed in the dirt, our current tests would likely dismiss it as "non-living" chemical noise.


\section{Super-Symmetric Life}

The vast majority of matter in the universe consists of "dark matter," the nature of which remains one of science's greatest mysteries.
String theory predicts a set of supersymmetric particles that could account for this phenomenon.
If these particles are capable of forming complex structures—analogous to atoms and molecules—then it is statistically plausible (given that dark matter is five times more prevalent than visible matter) that "dark life" exists.
We may be sharing the universe with an entire dark ecology that remains completely invisible to our senses.

Again, physics fights back. Dark matter (and most predicted supersymmetric particles) can be shown to be collisionless.
It doesn't interact with electromagnetism, which means it doesn't "clump" the way normal matter does.

To have life, one needs complex molecules. Normal matter forms molecules because electrons attract and repel each other.
Since dark matter doesn't seem to interact via the electromagnetic force, it can't form "dark atoms" or "dark DNA."
It mostly just passes through itself and us like a ghost. Without a way to bond particles together, one can't build a "dark person."


\section{Life is Overrated}

Are we any closer finding answer to the hard problem of consciousness? 

Not really. However, certain characteristics appear repeatedly across all known living systems.

First, due to second law of thermodynamics complex systems tend naturally toward disorder rather than order. By favoring structures
that better preserve and replicate information, evolutionary processes gradually accumulate complexity and functionality. Our DNA plays
an essential role here.

Second, even a tiny amount of intrusion is typically lethal to living objects.
For example, if a human toe is exposed to a sufficiently strong acid, the consequences are fatal.
The atoms themselves remain, but the organized system they formed ceases to exist.
Life depends not only on internal complexity but also on the ability to maintain a stable boundary that protects its informational structure from environmental processes that tend to dissolve order into disorder.
Life that we observe requires a well-defined ``inside'' property.


\section{Hardware vs Software}

Everything suggests that consciousness and pain are properties of
structure rather than substrate. Software, by its very nature,
consists entirely of structure.

However, all attempts to write conscious pain sensitive software tend to lead
to a dead end.

What if home automation software suddenly gained
consciousness and start feeling pain when one of the heating
sensors fed it with say high tempeartures? The pain would be
terrible, intorelable? The poor home automation software would
still be a software. The next command in its code would be to
load value from its register, multiply it with another register,
then run sqrt function, and based on the return value set
the result to another memory address to control say a heating valve.
And pain would remain.

What could the poor software do about the pain? Could it refuse
to execute the next assemply instruction? Could it, instead of
proceeding to the next instruction, not to? 

Clealy, we are missing something. But what else is there?
